% from ega0-20-fr.tex
\oldpage[0\textsubscript{IV-1}]{130}
En outre, comme le diagramme
\[
\xymatrix{
  \mathrm{P}^1_{C/A} \ar[r]^{v'} & \mathrm{P}^1_{C/B} \\
  C \ar[u]^{p_2} \ar[r]^{\mathrm{I}_C} & C \ar[u]_{p_2}
}
\tag{20.5.3.4}
\label{0.20.5.3.4-fr}
\]
est aussi commutatif, on en déduit que le diagramme
\[
\xymatrix{
  \Omega^1_{C/A} \ar[r]^{v_{C/B/A}} & \Omega^1_{C/B} \\
  C \ar[u]^{d_{C/A}} \ar[r]^{\mathrm{I}_C} & C \ar[u]_{d_{C/B}}
}
\tag{20.5.3.5}
\label{0.20.5.3.5-fr}
\]
est commutatif.

Enfin, si $s:A'\to A$ est un second homomorphisme d'anneaux, on a la
propriété de transitivité
\begin{equation}
\tag{20.5.3.6}
\label{0.20.5.3.6-fr}
  (v\circ s)_{C/B/A'}=v_{C/B/A}\circ s_{C/A/A'}.
\end{equation}
\end{env}

\begin{env}[20.5.4]
\label{0.20.5.4-fr}
Si maintenant on a un diagramme commutatif d'homomorphismes d'anneaux
\[
\xymatrix{
  B \ar[r]^v & B' \\
  A \ar[u]^f \ar[r]^u & A' \ar[u]_g
}
\]
on déduit de \sref{0.20.5.2.4-fr} et \sref{0.20.5.3.3-fr}, par composition,
un homomorphisme de $B'$-modules
\begin{equation}
\tag{20.5.4.1}
\label{0.20.5.4.1-fr}
  \Omega^1_{B/A}\otimes_B B'
  \xrightarrow{v_{B'/B/A}}\Omega^1_{B'/A}
  \xrightarrow{u_{B'/A'/A}}\Omega^1_{B'/A'}
\end{equation}
tel que le diagramme de $A'$-homomorphismes
\[
\xymatrix{
  \Omega^1_{B/A}\otimes_B B' \ar[r] & \Omega^1_{B'/A'} \\
  B' \ar[u]^{d_{B/A}\otimes1} \ar[r]^{\mathrm{I}_{B'}} &
  B' \ar[u]_{d_{B'/A'}}
}
\tag{20.5.4.2}
\label{0.20.5.4.2-fr}
\]
soit commutatif.

L'homomorphisme \sref{0.20.5.4.1-fr} correspond d'ailleurs à un
di-homomorphisme de $B$-modules
\begin{equation}
\tag{20.5.4.3}
\label{0.20.5.4.3-fr}
  \Omega^1_{B/A}\to\Omega^1_{B'/A'}.
\end{equation}
\end{env}

\begin{proposition}[20.5.5]
\label{0.20.5.5-fr}
Si $A'$, $B$ sont deux $A$-algèbres et $B'=B\otimes_A A'$,
l'homomorphisme canonique \sref{0.20.5.4.1-fr}
\begin{equation}
\tag{20.5.5.1}
\label{0.20.5.5.1-fr}
  \Omega^1_{B/A}\otimes_B B'\to\Omega^1_{B'/A'}
\end{equation}
est bijectif.
\end{proposition}

\begin{proof}
